\gddchapter{Testing sequence}
The prototype can be played with both voice and keyboard and the navigation modes can be swapped with a button press. 

\section{Baseline A - Keyboard navigation}
At first the participant will spend 10 minutes by playing the game using keyboard navigation.
This allows the player to anchor in the typical experience and serves as a comparative point in the A/B testing.

\subsection{Observation log}
\begin{table}[h]
\centering
\begin{tabular}{|p{5cm}|p{10cm}|}
\hline
\textbf{Field} & \textbf{Response} \\ \hline
\textbf{Physical comfort} How is the participant feeling & The player is very comfortable and navigating quickly \\[1.5cm] \hline
\textbf{Immersion markers} (flow \& frustration) & The player instantly went deep into the experience and seemed to be deep into it during the 10 minutes\\[1.5cm] \hline
\end{tabular}
\end{table}

\section{Novel version B - Voice navigation}
After 10 minutes the researcher switches the navigation to voice input mode by using the keyboard. The prototype is not restarted,
instead the player is instructed to continue his playthrough using voice. The decision not to restart is placed to reduce any frustration that the player
might feel if they have to restart and repeat the actions that they just did.

\subsection{Observation log}
\begin{table}[H]
\centering
\begin{tabular}{|p{5cm}|p{10cm}|}
\hline
\textbf{Field} & \textbf{Response} \\ \hline
\textbf{Physical comfort} How is the participant feeling & The player looks perfectly comfortable\\[1.5cm] \hline
\textbf{Natural language observation} (is the player using long sentences or simple commands? What type of language do they use? Are they speaking first or 3rd person?) & The player is using simple commands using the structure of "Let's X"\\[1.5cm] \hline
\textbf{Linguistic guessing} (is the particpant trying to guess the possible commands by speaking to the game?) & The participant seems to know the pre-exisitng limits from the keyboard version but he's also trying to explore the possiblities and limits of the voice driven interacion.\\[1.5cm] \hline
\textbf{Processing time user reaction} (Does the wait time break the suspension of disbelief)& There seems to be a small, yet noticable impact of the processing time on the player immersion.\\[1.5cm] \hline
\textbf{Voice processing pipeline edge cases} (System edge cases eg. two vs to)& The voice processing engine teleported the player across the map and failed to recognise his voice numerous times.\\[1.5cm] \hline
\end{tabular}
\end{table}

\subsection{Sticky moments}
The moments that stood out during testing marked with timestamps.
\begin{table}[h]
\centering
\begin{tabular}{|p{3cm}|p{10cm}|}
\hline
\textbf{Timestamp} & \textbf{Log} \\ \hline
04:43 & Start of the playtest \\ \hline 
08:10 & The fish store has a missing image \\ \hline 
11:05 & Stairway up \& down bug \\ \hline 
12:40 & Bug - can't pick up brother \\ \hline 
13:40 & Bug - the player can leave the building \\ \hline 
14:37 & The player expressed the want for more interaction \\ \hline 
23:00 & Bug - the leg stabiliser can't be picked up \\ \hline 
26:00 & Bug - empty pistol pickup \\ \hline 
26:30 & Bug - double teleportation \\ \hline 



\end{tabular}
\end{table}

\section{Alignment check}
How much was the researcher participating in the gameplay experience? 
Is the participant related in any way to the researcher? Could that influence the result?

The researcher did his best to not defend the system based on the experiences from the previous test. The participant's background is in game development which can have an influence on the test result, since people who make games look at them differently than just regular players.
